% Revize 2026-10-28: obsah, matematika a návaznost; podrobnosti a časy
% viz planovani/REVIZE_PRIPRAV_1_ROCNIK.md.
% Organizační hodina: v matematické knihovně nemá obsahový protějšek.
% 01_ÚVD_ÚvodníHodina.tex
% Ucitelska priprava: 1. rocnik, hodina 01, organizace vyuky matematiky.
% Kompilace z Work/MAT: latexmk -cd -pdf 1_rocnik/pripravy/01_ÚVD_ÚvodníHodina.tex
% Vyzaduje Work/MAT/sablony/pripravy/priprava.sty (verze 1.1).
\documentclass[11pt,a4paper]{article}
% Spolecny styl pro ucitelske pripravy MAT (Work/MAT)
% Verze 1.1, 2026-10-23
\NeedsTeXFormat{LaTeX2e}
\ProvidesPackage{priprava}[2026/10/23 v1.1 Sablona ucitelskych priprav MAT]

% Zakladni typografie a cestina
\RequirePackage[T1]{fontenc}
\RequirePackage[utf8]{inputenc}
\RequirePackage[czech]{babel}
\RequirePackage{lmodern}
\RequirePackage{microtype}
\RequirePackage{amsmath,amssymb}
\RequirePackage{xcolor}
\RequirePackage{enumitem}
\RequirePackage[most]{tcolorbox}
\RequirePackage{tabularx,booktabs,array}
\RequirePackage{fancyhdr}
\RequirePackage{needspace}
\RequirePackage{geometry}
\RequirePackage[hidelinks]{hyperref}

\geometry{a4paper,top=17mm,bottom=18mm,left=23mm,right=23mm}
\definecolor{notegray}{gray}{0.39}
\setlength{\parindent}{0pt}
\setlength{\parskip}{2.5pt plus 1pt minus 0.5pt}
\setlist[itemize]{topsep=2pt,itemsep=1pt,parsep=0pt,leftmargin=1.45em}
\setlist[enumerate]{topsep=2pt,itemsep=2pt,parsep=0pt,leftmargin=1.65em}
\setcounter{secnumdepth}{0}
\renewcommand{\arraystretch}{1.12}

% Zahlavi: vlevo rocnik | globalni tema, vpravo cislo: nazev hodiny
% Zapati zustava zcela prazdne.
\pagestyle{fancy}
\fancyhf{}
\renewcommand{\headrulewidth}{0.35pt}
\renewcommand{\footrulewidth}{0pt}
\setlength{\headheight}{14pt}
\newcommand{\nastavPripravu}[4]{%
  \fancyhead[L]{\small #1\enspace\textbar\enspace #2}%
  \fancyhead[R]{\small #3: #4}%
}

% Minimalisticke cernobile boxy. Kazda definice zvlast.
\tcbset{commonlesson/.style={enhanced,breakable=false,colback=white,colframe=black,
  colbacktitle=white,coltitle=black,boxrule=0.45pt,
  arc=0.8mm,outer arc=0.8mm,left=3mm,right=3mm,top=1.6mm,bottom=1.6mm,
  toptitle=1mm,bottomtitle=0.7mm,fonttitle=\bfseries,
  attach boxed title to top left={xshift=2.5mm,yshift=-1.15mm},
  boxed title style={colback=white,colframe=white,boxrule=0pt,left=1mm,right=1mm}}}
\newtcolorbox{zadaniBox}[1][]{commonlesson,title={Zadání},#1}
\newtcolorbox{prikladBox}[1][]{commonlesson,title={Příklad},#1}
\newtcolorbox{definiceBox}[1][]{commonlesson,title={Definice},#1}

% Spolecny box pro zadani a jeho reseni, bez mezery mezi castmi.
% Pouziti:
%   \begin{prikladReseniBox}
%     Zadani ulohy.
%     \tcblower
%     \begin{align*} ... \end{align*}
%   \end{prikladReseniBox}
% Spodni cast ma jemny sedy podklad a delici linku; lze v ni sazet
% vice radku matematiky. Obe casti maji jen jeden spolecny ramecek.
\tcbset{solutionlesson/.style={
  commonlesson,
  enhanced jigsaw,
  breakable,
  colbacklower=black!3,
  segmentation style={solid,draw=black!25,line width=0.35pt},
  fontlower=\small\color{notegray},
  before lower={\noindent\textbf{Řešení:}\enspace},
  leftlower=3mm,rightlower=3mm
}}
\newtcolorbox{prikladReseniBox}[1][]{solutionlesson,title={Příklad},#1}
\newtcolorbox{zadaniReseniBox}[1][]{solutionlesson,title={Zadání},#1}

% Pomucky pro autora pripravy: sedive poznamky; \reseni pro starsi pripravy
\newcommand{\poznamka}[1]{\par\begingroup\small\color{notegray}\itshape
  \textnormal{\textbf{Poznámka pro učitele.}}\enspace #1\par\endgroup}
\newcommand{\reseni}[1]{\par\begingroup\small\color{notegray}
  \textbf{Očekávané řešení:}\enspace #1\par\endgroup}
\newcommand{\blok}[1]{\needspace{5\baselineskip}\subsection*{#1}}
\newcommand{\vyrok}[1]{\uv{#1}}

\nastavPripravu{1. ročník}{Organizace výuky}{01}{Úvodní hodina}

\begin{document}

\blok{Před hodinou}
\begin{itemize}[label=$\square$]
  \item Otevřít příslušný kurz matematiky v \textbf{Moodle}; zkontrolovat, že jsou pravidla zveřejněná a aktuální.
  \item Připravit přístup k \textbf{Bakalářům} a ověřit seznam žáků.
  \item Připravit odkazy na materiály, tematický plán a případné úvodní dokumenty.
\end{itemize}

\blok{Seznámení s předmětem}
\begin{itemize}[label=$\square$]
  \item Stručně představit obsah matematiky v prvním ročníku a návaznost na další ročníky.
  \item \textbf{Projít se žáky přímo v Moodle podmínky klasifikace} a ukázat, kde je kdykoli najdou.
  \item Vysvětlit typy ověřování znalostí: průběžné testy, čtvrtletní a pololetní práce, případné ústní zkoušení a další formy hodnocení \emph{podle aktuálního znění pravidel}.
  \item Projít podmínky omlouvání, náhradních termínů, dopisování prací a případných oprav.
  \item Projít požadované \textbf{pomůcky} (sešit, rýsovací potřeby, kalkulačka podle pokynů) a pravidla práce s nimi.
  \item Připomenout pravidla pro \textbf{mobilní telefony, chytré hodinky} a jiné elektronické pomůcky zejména při písemných pracích.
  \item Ukázat, kde budou dostupné přípravy, pracovní listy, domácí úkoly a oznámení.
  \item Vysvětlit způsob komunikace a jak mají žáci postupovat při absenci nebo nejasnostech.
  \item Dát prostor pro dotazy; ověřit, že žáci vědí, kde přesné znění pravidel dohledat.
\end{itemize}

\poznamka{Konkrétní hranice známek, váhy, počty testů, náhradní termíny ani sankce neopisovat zpaměti. Vždy vycházet z aktuálních informací v Moodle a platných školních pravidel. Pokud se některý údaj změnil, opravit ho před jeho sdělením třídě.}

\blok{Zápis do Bakalářů}
\begin{zadaniBox}
\textbf{Téma hodiny:} Úvodní hodina -- organizace výuky a pravidla hodnocení v MAT.

\medskip
\textbf{Zápis:} Žáci byli seznámeni s pravidly práce a hodnocení v předmětu MAT, podmínkami klasifikace a organizací výuky.
\end{zadaniBox}

\poznamka{Formulaci použij až po skutečném seznámení žáků s pravidly; případné nepřítomné žáky informuj následně. Zápis uprav podle toho, co se v hodině opravdu uskutečnilo.}

\blok{Po hodině}
\begin{itemize}[label=$\square$]
  \item Zapsat téma a poučení žáků do Bakalářů.
  \item Zaznamenat dotazy a nejasnosti, které je třeba doplnit v Moodle nebo příště vysvětlit.
  \item Navázat hodinou \texttt{02\_VRL\_CoJeVyrok.tex}.
\end{itemize}

\end{document}
